\documentclass[12pt]{article}
\usepackage[a4paper,margin=1in]{geometry}
\usepackage{setspace}
\setstretch{1.15}

\begin{document}

\begin{center}
{\Large \textbf{Energy Generation:}}\\[4pt]
{\large \textbf{How Power Stations Work \& The Role of Fossil Fuels}}
\end{center}

Electricity generation is based on a simple idea: \textbf{you must turn a generator}. A generator produces electricity when a coil of wire spins inside a magnetic field. The challenge is finding a reliable way to keep that coil turning.

There are two main approaches:

\begin{itemize}
    \item \textbf{Direct generation}: The generator is turned directly by a natural force such as moving air or water. Examples include \textbf{wind turbines} and \textbf{hydroelectric dams}. These methods are often \textbf{renewable} and produce \textbf{no carbon dioxide during operation}. However, some are \textbf{unreliable} because the natural resource may vary (e.g.\ wind speed).
    
    \item \textbf{Indirect generation}: A fuel is burned or heat is released to produce \textbf{steam}, which then spins a turbine connected to the generator. This is how most traditional power stations work, including those using \textbf{coal, oil, gas, nuclear fuel}, and even \textbf{biomass}. Indirect generation can be highly \textbf{reliable}, because the fuel supply can be controlled.
\end{itemize}

When evaluating any method of electricity generation, four important concepts are often considered:

\begin{itemize}
    \item \textbf{Renewable}: Will the energy source run out?
    \item \textbf{Reliable}: Can electricity be generated consistently when needed?
    \item \textbf{CO$_2$ emissions}: Does the process release carbon dioxide, a greenhouse gas?
    \item \textbf{Airborne pollutants}: Does the method release substances such as sulphur dioxide, nitrogen oxides, or particulates that can harm health and the environment?
\end{itemize}

\section*{Fossil Fuel Power Stations}

Fossil fuel power stations burn \textbf{coal, oil}, or \textbf{natural gas} to heat water, producing steam that spins a turbine. The steam turns the generator, producing electricity. This method has been widely used for decades because it is straightforward, scalable, and efficient at producing large amounts of power.

\subsection*{Advantages}

\begin{itemize}
    \item Fossil fuels provide a \textbf{reliable} power supply because they can be stored and burned on demand.
    \item They allow for \textbf{continuous electricity generation}, regardless of weather or time of day.
    \item Existing global infrastructure means fossil fuel power stations are relatively cheap to build compared to newer technologies.
    \item Gas power stations in particular can change output quickly, making them useful for balancing the electricity grid.
\end{itemize}

\subsection*{Disadvantages}

\begin{itemize}
    \item Fossil fuels are \textbf{non-renewable} and will eventually run out.
    \item Burning fossil fuels produces large amounts of \textbf{CO$_2$}, contributing to climate change.
    \item They release \textbf{airborne pollutants} such as sulphur dioxide, nitrogen oxides, and particulates, which can cause acid rain and respiratory problems.
    \item Coal and oil power stations have high overall emissions and environmental impact.
    \item Mining and transporting fossil fuels can damage ecosystems and create additional pollution.
\end{itemize}

Fossil fuel power stations remain important in many countries due to their high reliability and output, but their environmental impact has driven a global shift toward renewable energy sources and cleaner technologies.

\section*{Tasks}

\begin{enumerate}
    \item Explain the difference between \textbf{direct} and \textbf{indirect} electricity generation. Include an example of each.
    
    \item Define the terms \textbf{renewable} and \textbf{reliable}. Give one example of an energy source that is renewable but unreliable, and one that is reliable but non-renewable.
    
    \item Describe in your own words why nearly all power stations, regardless of type, require a generator to be turned.
    
    \item Write a detailed paragraph explaining how a fossil fuel power station generates electricity, starting from burning the fuel and ending with electricity being produced.
    
    \item Compare fossil fuel power stations with one renewable method of your choice. Explain how they differ in environmental impact and reliability.
    
    \item Suggest one reason why countries continue to use fossil fuels even though they produce pollution and CO$_2$. Use ideas from the worksheet to support your answer.
\end{enumerate}

\end{document}
